\documentclass[12pt]{article}
\usepackage{amsmath,amssymb,amsthm}
\usepackage{mathtools}
\usepackage{geometry}
\geometry{margin=1in}

% Theorem environments
\newtheorem{definition}{Definition}[section]
\newtheorem{lemma}[definition]{Lemma}
\newtheorem{theorem}[definition]{Theorem}
\newtheorem{corollary}[definition]{Corollary}
\newtheorem{proposition}[definition]{Proposition}
\newtheorem{remark}[definition]{Remark}

% Custom commands
\newcommand{\B}{\mathbb{B}}
\newcommand{\C}{\mathbb{C}}
\newcommand{\R}{\mathbb{R}}
\renewcommand{\H}{\mathbb{H}}
\newcommand{\Tr}{\mathrm{Tr}}
\newcommand{\re}{\mathrm{Re}}
\newcommand{\im}{\mathrm{Im}}

\title{Step 4: Properties of the UBT Stress-Energy Tensor $T_{\mu\nu}$\\
\large Realness, Symmetry, and Divergence-Free Condition}
\author{UBT Theory Development}
\date{2026}

\begin{document}

\maketitle

\begin{abstract}
For the Einstein equations~$G_{\mu\nu} = 8\pi G\,T_{\mu\nu}$ to be
self-consistent, the stress-energy tensor $T_{\mu\nu}$ must satisfy three
properties: (1)~it must be real-valued; (2)~it must be symmetric
$T_{\mu\nu} = T_{\nu\mu}$; (3)~it must be divergence-free
$\nabla^\mu T_{\mu\nu} = 0$.  We prove each property, specifying exactly
what conditions on $\Theta$ are required and which hold identically from
the structure of the Hilbert prescription.
\end{abstract}

\tableofcontents

% ------------------------------------------------------------
\section{Property 1: Realness}

\begin{theorem}[Realness of $T_{\mu\nu}$]
\label{thm:real}
The physical stress-energy tensor $T_{\mu\nu} = \re(\mathcal{T}_{\mu\nu})$
is real-valued by definition:
\begin{equation}
  T_{\mu\nu} \in \mathbb{R} \quad \text{for all } \mu,\nu.
\end{equation}
\end{theorem}

\begin{proof}
  By Definition~4.1 of \texttt{step2\_biquat\_stress\_energy.tex},
  $T_{\mu\nu}$ is defined as the real part of the biquaternionic tensor
  $\mathcal{T}_{\mu\nu}$.  Since the real-part map
  $\re: \mathbb{B} \to \mathbb{R}$ projects onto the scalar (real)
  component of any biquaternion $b = b_0 + ib_1\mathbf{i} + b_2\mathbf{j}
  + b_3\mathbf{k}$ with $b_0,\ldots,b_3 \in \mathbb{R}$, the output is
  real by construction.
\end{proof}

\begin{remark}
  Realness is not a condition on $\Theta$; it follows automatically from
  the projection step $T_{\mu\nu} := \re(\mathcal{T}_{\mu\nu})$.
  The imaginary components of $\mathcal{T}_{\mu\nu}$ carry information
  about phase curvature in the UBT extended sector but do not enter
  the observable Einstein equations.
\end{remark}

% ------------------------------------------------------------
\section{Property 2: Symmetry $T_{\mu\nu} = T_{\nu\mu}$}

\begin{theorem}[Symmetry from Hilbert prescription]
\label{thm:symmetry}
The Hilbert stress-energy tensor is automatically symmetric:
\begin{equation}
  T_{\mu\nu} = T_{\nu\mu}.
\end{equation}
\end{theorem}

\begin{proof}
  The metric $g^{\mu\nu}$ is symmetric under index interchange:
  $g^{\mu\nu} = g^{\nu\mu}$.  The functional derivative
  $\delta S/\delta g^{\mu\nu}$ is taken with respect to a symmetric tensor,
  so the result must also be symmetric.  More precisely:
  \begin{equation}
    \frac{\delta S}{\delta g^{\mu\nu}} = \frac{\delta S}{\delta g^{\nu\mu}},
  \end{equation}
  because $g^{\mu\nu}$ and $g^{\nu\mu}$ are not independent components.
  Under the symmetrised variation $\delta g^{\mu\nu} = \delta g^{\nu\mu}$,
  the functional derivative is defined as the symmetric part:
  \begin{equation}
    \frac{\delta S}{\delta g^{(\mu\nu)}}
    = \tfrac{1}{2}\Bigl(\frac{\delta S}{\delta g^{\mu\nu}}
      + \frac{\delta S}{\delta g^{\nu\mu}}\Bigr),
  \end{equation}
  which is manifestly symmetric in $(\mu,\nu)$.
  Hence $T_{\mu\nu} = -(2/\sqrt{-g})\,\delta S/\delta g^{(\mu\nu)}
  = T_{\nu\mu}$.

  \textbf{Explicit verification from formulas.}
  From \texttt{step1\_metric\_variation.tex}:
  \begin{itemize}
    \item $T^{(\mathrm{kin})}_{\mu\nu}$: symmetric because
          $\Tr[(\nabla_\mu\Theta)^\dagger\nabla_\nu\Theta]
          = \Tr[(\nabla_\nu\Theta)^\dagger\nabla_\mu\Theta]^*$
          and the real part is symmetric in $(\mu,\nu)$; and
          $g_{\mu\nu}$ is symmetric.
    \item $T^{(\mathrm{pot})}_{\mu\nu} = -g_{\mu\nu}V$:
          symmetric because $g_{\mu\nu}$ is symmetric.
    \item $T^{(\mathrm{gauge})}_{\mu\nu}$: symmetric because
          $\Tr[F_{\mu\alpha}F_\nu{}^\alpha] = \Tr[F_{\nu\alpha}F_\mu{}^\alpha]$
          (cyclicity of trace and $F_{\mu\nu} = -F_{\nu\mu}$)
          and $g_{\mu\nu}$ is symmetric.
  \end{itemize}
  Each piece is symmetric in $(\mu,\nu)$, so $T_{\mu\nu} = T_{\nu\mu}$.
\end{proof}

\begin{remark}[Einstein tensor consistency]
  The Einstein tensor $G_{\mu\nu} = R_{\mu\nu} - \tfrac{1}{2}g_{\mu\nu}R$
  is also symmetric: $G_{\mu\nu} = G_{\nu\mu}$.  Theorem~\ref{thm:symmetry}
  confirms that both sides of $G_{\mu\nu} = 8\pi G\,T_{\mu\nu}$ have the
  same symmetry, as required.
\end{remark}

% ------------------------------------------------------------
\section{Property 3: Divergence-Free Condition}

\begin{theorem}[Divergence-free: Bianchi route]
\label{thm:divergence_bianchi}
If the metric $g_{\mu\nu}$ and field $\Theta$ jointly satisfy the full
system of equations (Einstein equations + matter field equations), then:
\begin{equation}
  \label{eq:conservation}
  \nabla^\mu T_{\mu\nu} = 0.
\end{equation}
This is \textbf{automatic} from the contracted Bianchi identity.
No additional condition on $\Theta$ is required.
\end{theorem}

\begin{proof}
  The contracted Bianchi identity states:
  \begin{equation}
    \nabla^\mu G_{\mu\nu} = 0
  \end{equation}
  identically for any smooth metric $g_{\mu\nu}$, independent of any
  field equation.  From Theorem~3.1 of
  \texttt{step3\_einstein\_with\_matter.tex},
  $G_{\mu\nu} = 8\pi G\,T_{\mu\nu}$.
  Applying $\nabla^\mu$ to both sides:
  \begin{equation}
    \nabla^\mu G_{\mu\nu} = 8\pi G\,\nabla^\mu T_{\mu\nu} = 0.
  \end{equation}
  Since $8\pi G \neq 0$, equation~\eqref{eq:conservation} follows.
\end{proof}

\begin{theorem}[Divergence-free: Noether route]
\label{thm:divergence_noether}
When $\Theta$ satisfies the Euler--Lagrange equations
$\nabla^\mu\nabla_\mu\Theta = \partial V/\partial\Theta^\dagger$
(Theorem~4.1 of \texttt{appendix\_AA\_theta\_action.tex}), the
diffeomorphism invariance of $S_{\mathrm{matter}}$ implies:
\begin{equation}
  \nabla^\mu T_{\mu\nu} = 0
\end{equation}
as a Noether identity.
\end{theorem}

\begin{proof}
  Under an infinitesimal diffeomorphism generated by $\xi^\mu$,
  the metric variation is $\delta g_{\mu\nu} = -\nabla_\mu\xi_\nu
  - \nabla_\nu\xi_\mu$ (Lie derivative) and the field transforms as
  $\delta\Theta = \xi^\mu\nabla_\mu\Theta$.
  Invariance of $S_{\mathrm{matter}}$ under this combined variation gives,
  after using the $\Theta$ equation of motion to eliminate $\delta S/\delta\Theta$:
  \begin{equation}
    0 = -\int d^4x\,\sqrt{-g}\,T^{\mu\nu}\nabla_\mu\xi_\nu
      = \int d^4x\,\sqrt{-g}\,(\nabla^\mu T_{\mu\nu})\xi^\nu.
  \end{equation}
  Since $\xi^\nu$ is arbitrary, $\nabla^\mu T_{\mu\nu} = 0$.
\end{proof}

\begin{remark}[Summary of conditions]
\label{rem:conditions}
\begin{center}
\begin{tabular}{|l|l|l|}
\hline
\textbf{Route} & \textbf{Condition on $\Theta$} & \textbf{Status} \\
\hline
Bianchi identity
  & None (follows from geometry alone)
  & \textbf{Automatic} \\
\hline
Noether (matter sector)
  & $\Theta$ on-shell (E-L equations satisfied)
  & \textbf{Automatic on-shell} \\
\hline
\end{tabular}
\end{center}
In either case, no additional assumptions are needed beyond the
existing UBT action principle.
\end{remark}

% ------------------------------------------------------------
\section{Summary Table}

\begin{theorem}[Complete property verification]
\label{thm:summary}
The UBT stress-energy tensor $T_{\mu\nu}$, derived from the Hilbert
prescription applied to $S_{\mathrm{matter}}[\Theta, g]$, satisfies:
\end{theorem}

\begin{center}
\renewcommand{\arraystretch}{1.4}
\begin{tabular}{|l|l|l|}
\hline
\textbf{Property} & \textbf{Result} & \textbf{Proof} \\
\hline
Realness: $T_{\mu\nu} \in \mathbb{R}$
  & $\checkmark$~Proved
  & Definition (real-part projection) \\
\hline
Symmetry: $T_{\mu\nu} = T_{\nu\mu}$
  & $\checkmark$~Proved
  & Hilbert prescription on symmetric $g^{\mu\nu}$ \\
\hline
Divergence-free: $\nabla^\mu T_{\mu\nu} = 0$
  & $\checkmark$~Proved
  & Bianchi identity + Noether \\
\hline
QED limit: $T_{\mu\nu} \to T^{(\mathrm{EM})}_{\mu\nu} + T^{(\mathrm{Dirac})}_{\mu\nu}$
  & $\checkmark$~Proved
  & Step~2, Proposition~5.1 \\
\hline
\end{tabular}
\end{center}

\begin{remark}[Outstanding items -- CANDIDATE]
\label{rem:candidate}
The following items are stated as candidates pending further work:
\begin{itemize}
  \item \textbf{CANDIDATE:} Whether $T_{\mu\nu}$ satisfies the Weak Energy
        Condition (WEC: $T_{\mu\nu}u^\mu u^\nu \geq 0$ for timelike $u^\mu$)
        depends on the specific field configuration $\Theta$ and is not
        proved in full generality.
  \item \textbf{CANDIDATE:} Conservation $\nabla^\mu T_{\mu\nu} = 0$ for
        the imaginary sector $\im(\mathcal{T}_{\mu\nu})$ is not addressed
        here; it requires separate analysis within the biquaternionic
        field equations.
\end{itemize}
\end{remark}

\begin{thebibliography}{99}
\bibitem{wald}    Wald, R.~M. (1984). \textit{General Relativity}.
                  University of Chicago Press.
\bibitem{step1}   UBT Theory Development.
                  \texttt{step1\_metric\_variation.tex}, 2026.
\bibitem{step2}   UBT Theory Development.
                  \texttt{step2\_biquat\_stress\_energy.tex}, 2026.
\bibitem{step3}   UBT Theory Development.
                  \texttt{step3\_einstein\_with\_matter.tex}, 2026.
\bibitem{appendix_AA}
                  UBT Theory Development.
                  \texttt{appendix\_AA\_theta\_action.tex}, 2025.
\end{thebibliography}

\end{document}
